\begin{solapartat}
\punts{0,25} Ens diuen que el nucli inicial té un excés de neutrons; per tant, s'estabilitzarà convertint un neutró en protó. Com que tindrem un protó de més, $Z$ passa d'11 a 12 i, per tant, el nou element serà el magnesi-24.

\punts{1,0} En la reacció de transformació d'un neutró en un protó, s'emet un electró per respectar la conservació de la càrrega elèctrica.

O alternativament, ${}^{1}_{0}n \to {}^{1}_{1}p + {}^{\;0}_{-1}e + {}^{0}_{0}\bar{\nu}_e$.

Per tant, la reacció serà:
\[ {}^{24}_{11}Na \to {}^{24}_{12}Mg + {}^{\;0}_{-1}e + {}^{0}_{0}\bar{\nu}_e \]
\end{solapartat}

\begin{solapartat}
\punts{0,75} Per trobar l'activitat al cap de 10~h, necessitem el coeficient de desintegració $\lambda$ que obtenim del temps de semidesintegració, que és el temps en reduir l'activitat a la meitat: $A(t) = \frac{A_0}{2} = A_0\,e^{-\lambda t_{1/2}}$ i per tant, $t_{1/2} = \frac{\ln 2}{\lambda}$.

Directament obtenim:
\[ \lambda = \frac{\ln 2}{t_{1/2}} = \frac{\ln 2}{15} = 4,62\cdot 10^{-2}\ \mathrm{h^{-1}} \]
L'activitat després de 10 hores:
\[ A = A_0\,e^{-\lambda t} = 5,80\cdot 10^{5}\,e^{-4,62\cdot 10^{-2}\cdot 10} = 3,65\cdot 10^{5}\ \mathrm{Bq} \]

\punts{0,5} Després de 10~h, l'isòtop de sodi injectat s'ha dissolt en tot el volum de sang. Sabent que l'activitat de 1,00~mL són 60~Bq i que la de tot el volum de sang és l'activitat que acabem de calcular, podem calcular el volum total de sang, ja que la proporció serà la mateixa:
\[ \frac{3,65\cdot 10^{5}\ \mathrm{Bq}}{V_{total}} = \frac{60\ \mathrm{Bq}}{1\ \mathrm{mL}} \]
Per tant:
\[ V_{total} = 3,65\cdot 10^{5}\ \mathrm{Bq}\cdot\frac{1\ \mathrm{mL}}{60\ \mathrm{Bq}}\,\frac{1\ \mathrm{L}}{1000\ \mathrm{mL}} = 6,09\ \mathrm{L} \]
\end{solapartat}
